% !PRINTED 597
\weberpage{597}
%<<CONT
in $\Omega$, then $Q$ is the group of the same equation in $\Omega(\varepsilon)$, or,
what is the same thing, in $\Omega(\varepsilon_1), \Omega(\varepsilon_2) \ldots \Omega(\varepsilon_{m-1})$. But now we can
apply the concluding theorem in \S.\,157. By this theorem the
index $j$ of the divisor $Q$ of $P$ must be a divisor of $m$, and
since $m$ is assumed to be a prime, we have $m = j$. Moreover, by
the same theorem, $\varepsilon$ is rationally representable by the roots of the
equation $f(x) = 0$:

\[
\varepsilon = \psi(x_0, x_1 \ldots x_{m-1}).
\]

This function belongs to the group $Q$, and if we apply to it
all the permutations of the group $P$, we obtain
the functions $\varepsilon, \varepsilon_1, \varepsilon_2 \ldots \varepsilon_{m-1}$ and no others. These
functions belong to the conjugate groups $\pi^{-1}Q\pi$. But since
each of these functions is rationally expressible by every other,
they must all belong to the same group,
that is, $Q$ is a normal divisor of $P$.

We have thus proved the first theorem:

\begin{quote}
\noindent I.\ \gesperrt{If the group $P$ of an equation is reduced by
the adjunction of the roots of an Abelian equation, then $P$ has a normal divisor $Q$ of prime index.}
\end{quote}

This theorem can also be reversed.

For if the group $P$ has a normal divisor $Q$ of
index $m$, then we can choose a function $\psi$ belonging to $Q$,
and the functions conjugate with it, $\psi, \psi_1, \psi_2 \ldots \psi_{m-1}$,
all belong to the same group. The field $\Omega(\psi)$ is a
normal field and $\psi$ the root of a normal equation. In the
field $\Omega(\psi)$, $Q$ is the group of $f(x) = 0$ (\S.\,155). But if
$m$ is a prime, then by \S.\,163 $\psi$ is the root of a
cyclic equation, and thus it is proved:

\begin{quote}
\noindent II.\ \gesperrt{If the group $P$ of $f(x) = 0$ possesses a normal divisor $Q$ of prime index $m$, then by
the adjunction of the root of a cyclic
equation of degree $m$ the group $P$ is
reduced to $Q$.}
\end{quote}

\begin{center}
\S.\,176.

\gesperrt{Metacyclic Equations.}
\end{center}

We shall call an equation whose complete solution can be
reduced to a chain of cyclic equations
%CONT>>
